\documentclass[11pt]{article}
\usepackage[a4paper,margin=1in]{geometry}
\usepackage{amsmath,amssymb,amsthm,mathtools}
\usepackage{graphicx}
\usepackage[hidelinks]{hyperref}
\emergencystretch=3em

\newtheorem{theorem}{Theorem}
\newtheorem{lemma}[theorem]{Lemma}
\newcommand{\Ga}{\Gamma}

\title{A Three-State Counterexample:\\
Infinitesimal Optimality Does Not Imply\\
Global Markovian-Coupling Optimality}
\author{A full counterexample to Remark 2.20 of arXiv:2509.23086}
\date{August 11, 2026}

\begin{document}
\maketitle

\begin{abstract}
Kang and Lim ask whether a coupling generator that minimizes the instantaneous
cost derivative at every state must minimize the expected cost at every later
time.  We give a counterexample on the three-point discrete space.  Both
marginal chains have deterministic unit-rate jumps, the cost has nonnegative
integer entries, and the instantaneous optimizer is unique at all nine joint
states.  Nevertheless, at time $t=2$ it is strictly worse than a coupling that
differs at only one state.  Uniformization gives the exact gap
$2e^{-4}(e^2-5)>0$.
\end{abstract}

\noindent\textbf{Status.}
\texttt{full counterexample; likely valid}.  Human review is recommended.

\section{The source question}

Let $\mathcal A,\mathcal B$ be Feller generators on a state space $\Pi$, let
$c:\Pi^2\to[0,\infty)$, and write $\Ga(\mathcal A,\mathcal B)$ for their
Feller coupling generators.  The source calls $\mathcal J_*$ $c$-optimal when
it attains, simultaneously at every $(x,y)$,
\[
 (\mathcal J_*c)(x,y)
 =\inf_{\mathcal J\in\Ga(\mathcal A,\mathcal B)}(\mathcal Jc)(x,y).
\]
Global-in-time optimality is the stronger condition
\begin{equation}\label{eq:global}
 (e^{t\mathcal J_*}c)(x,y)\le(e^{t\mathcal J}c)(x,y)
 \quad(t\ge0,\ (x,y)\in\Pi^2,\ \mathcal J\in\Ga(\mathcal A,\mathcal B)).
\end{equation}
Proposition 2.19 proves that \eqref{eq:global} implies $c$-optimality;
Remark 2.20 requests a counterexample to the converse.

\begin{center}
\includegraphics[width=.96\linewidth]{figures/open_question_crop.png}
\end{center}

\section{Construction}

Take $\Pi=\{0,1,2\}$ with the discrete topology, and define
\[
 f=(1,2,0),\qquad g=(1,0,1).
\]
Thus the marginal generators are
\[
 (\mathcal A\varphi)(x)=\varphi(f(x))-\varphi(x),\qquad
 (\mathcal B\psi)(y)=\psi(g(y))-\psi(y).
\]
Use the cost matrix, with rows indexed by $x$ and columns by $y$,
\begin{equation}\label{eq:cost}
 C=(c(x,y))_{x,y=0}^2=
 \begin{pmatrix}7&6&0\\10&2&10\\6&9&1\end{pmatrix}.
\end{equation}

Every coupling generator is determined at each $(x,y)$ by one number
$s_{xy}\in[0,1]$.  Its three possible transition rates are
\[
 q((x,y),(f(x),g(y)))=s_{xy},\quad
 q((x,y),(f(x),y))=q((x,y),(x,g(y)))=1-s_{xy}.
\]
Indeed, the two marginal rate constraints force exactly this form.  Conversely,
every such choice defines a coupling generator.  Since the space is finite and
discrete, all these generators are Feller and $c$ lies in all their domains.

Define $\mathcal J_*$ by
\begin{equation}\label{eq:star}
 (s^*_{xy})=
 \begin{pmatrix}1&0&1\\0&1&0\\1&0&1\end{pmatrix}.
\end{equation}
Thus $s=1$ means a simultaneous jump, while $s=0$ means that the two possible
one-coordinate jumps each occur at rate one.

\section{The generator is the unique instantaneous optimizer}

Put
\[
 D_{xy}=c(f(x),g(y))-c(f(x),y)-c(x,g(y))+c(x,y).
\]
For fixed $(x,y)$, the value $(\mathcal Jc)(x,y)$ is an affine function of
$s_{xy}$ whose coefficient is $D_{xy}$.  Direct substitution in
\eqref{eq:cost} gives
\begin{equation}\label{eq:D}
 (D_{xy})=
 \begin{pmatrix}-7&7&-14\\11&-11&16\\-4&4&-2\end{pmatrix}.
\end{equation}
Every entry is nonzero.  The unique minimizer is therefore $s_{xy}=1$ when
$D_{xy}<0$ and $s_{xy}=0$ when $D_{xy}>0$, which is precisely
\eqref{eq:star}.  Hence $\mathcal J_*$ is the unique $c$-optimal coupling
generator in the source's infinitesimal sense.

Now let $\mathcal L$ agree with $\mathcal J_*$ at every joint state except
\[
 s^{\mathcal L}_{22}=0
 \quad\text{instead of}\quad s^*_{22}=1.
\]
We prove that $\mathcal L$ beats $\mathcal J_*$ at $(2,2)$ and time $2$.

\section{Exact failure at time two}

Uniformize at rate two.  Put
\[
 R_*=2I+\mathcal J_*,\qquad R_{\mathcal L}=2I+\mathcal L.
\]
Both $R_*/2$ and $R_{\mathcal L}/2$ are stochastic matrices.  For
$S=(2,2)$ set
\[
 \Delta_k=(R_*^k c)(S)-(R_{\mathcal L}^k c)(S).
\]

We record a short exact recurrence.  Label
\[
 A=01,\ B=11,\ C=00,\ D=20,\ E=02,\ F=21.
\]
These six states never reach $S$, and the two matrices agree on them.  If
$X_k=(R^kc)(X)$ for one of these common states, then
\begin{align*}
 A_{k+1}&=B_k+C_k,& B_{k+1}&=B_k+D_k,&
 C_{k+1}&=C_k+B_k,\\
 D_{k+1}&=D_k+A_k,& E_{k+1}&=E_k+B_k,&
 F_{k+1}&=A_k+D_k.
\end{align*}
The initial values are
\[
 (A_0,B_0,C_0,D_0,E_0,F_0)=(6,2,7,6,0,9).
\]
Under $R_*$ the $S$-row is $v_S+v_A$; under $R_{\mathcal L}$ it is
$v_E+v_F$.  Thus a direct calculation gives
\[
 \Delta_0=0,\qquad \Delta_1=-2,\qquad \Delta_2=2.
\]
For completeness, the displayed recurrences imply $A_k=C_k$ for $k\ge1$ and
\[
 F_k-F_{k-1}=A_{k-1}=C_{k-1}\qquad(k\ge2).
\]
Together with $E_k-E_{k-1}=B_{k-1}$, this shows that from $k=2$ onward the
two $S$-iterates have the same increment $A_k$.  Hence
\begin{equation}\label{eq:deltas}
 \Delta_k=2\qquad(k\ge2).
\end{equation}

The uniformization series is absolutely convergent and gives
\begin{align*}
 \bigl((e^{2\mathcal J_*}-e^{2\mathcal L})c\bigr)(S)
 &=e^{-4}\sum_{k=0}^{\infty}\frac{2^k}{k!}\Delta_k\\
 &=e^{-4}\left(-4+2\sum_{k=2}^{\infty}\frac{2^k}{k!}\right)\\
 &=2e^{-4}(e^2-5)>0.
\end{align*}
The last inequality follows already from the strict Taylor bound
$e^2>1+2+2^2/2=5$.  Thus
\[
 (e^{2\mathcal J_*}c)(2,2)>(e^{2\mathcal L}c)(2,2),
\]
contradicting \eqref{eq:global}.

\begin{theorem}
Infinitesimal $c$-optimality of a Feller coupling generator does not, in
general, imply global-in-time $c$-optimality.  This remains false even on a
three-point state space, with deterministic unit-rate marginals, a
nonnegative integer cost, and a unique instantaneous optimizer.
\end{theorem}

\section{Verification and scope}

The accompanying script recomputes \eqref{eq:D}, enumerates all $2^9$ extreme
coupling policies to audit pointwise minimality, and checks the integer
recurrence \eqref{eq:deltas}.  No floating-point computation is used in the
proof.  The counterexample concerns a general finite-state cost and therefore
does not contradict the source's positive all-time result for translation-
invariant L\'evy generators under quadratic Euclidean cost.

A bounded search of the exact question, its distinctive wording, the arXiv
identifier, and later coupling literature found special simultaneous-
optimality theorems but no matching explicit counterexample.  Human
mathematical and literature review is nevertheless recommended.

\begin{thebibliography}{9}
\bibitem{KangLim}
W. Y. Kang and T. S. Lim,
\emph{On Optimal Markovian Couplings of L\'evy Processes},
arXiv:2509.23086 (2025), Remark 2.20.
\url{https://arxiv.org/abs/2509.23086}

\bibitem{KMM}
W. S. Kendall, M. B. Majka, and A. Mijatovi\'c,
\emph{Optimal Markovian coupling for finite activity L\'evy processes},
Bernoulli 30 (2024), 2821--2845.
\end{thebibliography}

\end{document}

